%!TeX root = ../main.tex


% ================================================================================================
% 4. Problemas de emparelhamento em grafos bipartidos
% ================================================================================================
\chapter{Problemas de emparelhamento em grafos bipartidos}
\label{cap:bipartidos}

Este capítulo apresenta os quatro problemas de emparelhamento em grafos bipartidos estudados neste trabalho. Todos compartilham a mesma estrutura de entrada, um grafo cujos vértices se dividem em dois conjuntos disjuntos, com arestas apenas entre eles, e diferem quanto ao objetivo da otimização:

\begin{alineas}
    \item \textbf{emparelhamento de cardinalidade máxima:} maximizar o número de arestas do emparelhamento, sem pesos;
    \item \textbf{problema de atribuição:} encontrar o emparelhamento perfeito de menor custo total;
    \item \textbf{problema de atribuição gargalo:} encontrar o emparelhamento perfeito cuja aresta mais cara seja a mais barata possível;
    \item \textbf{emparelhamento estável:} encontrar um emparelhamento em que nenhum par prefira estar junto a permanecer com os parceiros atuais, a partir de listas de preferências.
\end{alineas}

Cada problema é apresentado na mesma sequência: a definição e as propriedades necessárias aos algoritmos, os algoritmos exatos, os algoritmos aproximados e as heurísticas, classificados conforme a Seção~\ref{sec:tipos_metodo}, e uma tabela-resumo com a complexidade, o tipo e uma breve descrição de cada método.

% ------------------------------------------------------------------------------------------------
% 4.1. Emparelhamento de cardinalidade máxima
% ------------------------------------------------------------------------------------------------
\subsection{Emparelhamento de cardinalidade máxima}
\label{sec:emp_card_max_bip}

\subsubsection{Descrição do problema}

Dado um grafo não-direcionado $G=(V,E)$, o problema de emparelhamento de cardinalidade máxima (MCM) busca encontrar um emparelhamento $M \subseteq E$ tal que o número de arestas em $M$ seja maximizado, ou seja, $|M|$ é o maior entre todos os emparelhamentos possíveis em $G$.

Um caso especial, mas importante, do MCM é o \textbf{Emparelhamento perfeito}, que é um emparelhamento onde todos os vértices do grafo são cobertos por exatamente uma aresta do emparelhamento. Para que um emparelhamento perfeito exista, o grafo deve ter um número par de vértices e o emparelhamento máximo deve ter cardinalidade igual a $|V|/2$ \cite{jungnickel}.

\subsubsection{Propriedades}

Antes de continuar para os algoritmos e soluções, é necessário entender algumas propriedades importantes

\paragraph{Emparelhamento máximo vs maximal}
\label{p:emp_max_maximal}

Um emparelhamento máximo é aquele que contém o maior número possível de arestas, enquanto um emparelhamento maximal é aquele que não pode ser aumentado adicionando mais arestas, mas não necessariamente é o maior possível \cite{manber}.

\paragraph{Caminhos aumentantes}
\label{p:caminhos_aumentantes}

Um caminho aumentante $P$ em um grafo $G$ é chamado de caminho aumentante se $P$ tem tamanho ímpar, começa e termina em vértices livres (não cobertos por nenhuma aresta do emparelhamento atual) e alterna entre arestas que pertencem a $M$ e arestas que não pertencem a $M$. A existência de um caminho aumentante indica que o emparelhamento atual não é máximo \cite{schrijver}

\paragraph{Teorema de Kőnig}
\label{p:teorema_konig}

O teorema de Kőnig estabelece uma relação fundamental entre emparelhamentos e coberturas em grafos bipartidos. Ele afirma que, em qualquer grafo bipartido, o tamanho do maior emparelhamento é igual ao tamanho da menor cobertura de vértices \cite{jungnickel}.

A cobertura de vértices é um conjunto de vértices tal que cada aresta do grafo é incidente a pelo menos um vértice desse conjunto. O teorema de Kőnig é uma ferramenta poderosa para resolver problemas de emparelhamento, utilizado nos algoritmos de Hopcroft-Karp e no método Húngaro. A seguir será detalhado um método algorítmico para encontrar a cobertura de vértices mínima a partir de um emparelhamento máximo.

O método é definido pelos seguintes passos \cite{jungnickel}:
\begin{enumerate}
    \item Identifica-se o conjunto de vértices livres (não saturados por $M$) na partição esquerda $U$.
    \item A partir desses vértices, inicia-se uma travessia no grafo (utilizando Busca em Largura ou Profundidade) com regras estritas de alternância: transições do conjunto $U$ para $V$ devem ocorrer exclusivamente através de arestas que \textbf{não} pertencem ao emparelhamento ($E \setminus M$), enquanto as transições de $V$ de volta para $U$ devem utilizar obrigatoriamente as arestas do emparelhamento ($M$).
    \item Registra-se em um conjunto $Z$ todos os vértices que foram alcançados durante essa busca, tanto da partição $U$ quanto da partição $V$.
    \item A Cobertura Mínima de Vértices $C$ é construída selecionando-se os vértices de $U$ que \textbf{não} foram visitados e os vértices de $V$ que \textbf{foram} visitados. Matematicamente, a cobertura é dada por $C = (U \setminus Z) \cup (V \cap Z)$.
\end{enumerate}

Essa extração algorítmica possui complexidade linear em relação ao tamanho do grafo, $O(V + E)$, visto que requer apenas uma travessia simples após o cálculo do emparelhamento. Sua implementação é uma sub-rotina fundamental para métodos exatos de problemas mais complexos, como a atualização de limiares no Método Dual para o problema de atribuição gargalo \cite{cormen}.

\paragraph{Teorema de Hall}

O Teorema de Hall, frequentemente referido como Teorema do Casamento de Hall, é outro resultado estrutural para grafos bipartidos. Enquanto o teorema de König determina a cardinalidade do emparelhamento máximo, o teorema de Hall estabelece a condição necessária e suficiente para a existência de um emparelhamento que cubra integralmente um dos subconjuntos da partição \cite{schrijver}. 

Para um grafo bipartido $G=(U \cup V, E)$, o teorema afirma que existe um emparelhamento cobrindo todo o conjunto $U$ se, e somente se, para qualquer subconjunto $S \subseteq U$, o tamanho de sua vizinhança $N(S)$ (o conjunto de vértices em $V$ adjacentes a pelo menos um vértice em $S$) for maior ou igual ao tamanho de $S$. Matematicamente, exige-se que $|N(S)| \ge |S|$ para todo $S \subseteq U$ \cite{jungnickel}.  Essa propriedade, conhecida como Condição de Hall, é essencial para provar matematicamente a viabilidade de atribuições e identificar gargalos estruturais que impedem a formação de um emparelhamento perfeito, relacionando-se intimamente com os conceitos de cortes mínimos em redes de fluxo.

\paragraph{Teorema de Gallai e Conjunto Independente Máximo}

Um conjunto independente em um grafo é um subconjunto de vértices no qual não existem duas extremidades adjacentes, ou seja, nenhuma aresta conecta dois vértices desse conjunto. O problema de encontrar o Conjunto Independente Máximo (\textit{Maximum Independent Set} - MIS) é classicamente conhecido por pertencer à classe NP-Difícil para grafos gerais \cite{garey-jhonson}. No entanto, quando o grafo é restrito à classe dos grafos bipartidos, o problema pode ser resolvido em tempo polinomial graças à sua profunda relação com os emparelhamentos e coberturas.

A base dessa resolução é o Teorema de Gallai, proposto em 1959. O teorema estabelece que, para qualquer grafo $G=(V, E)$ que não possua vértices isolados, a soma do tamanho da Cobertura Mínima de Vértices (MVC) com o tamanho do Conjunto Independente Máximo (MIS) é sempre igual ao número total de vértices do grafo \cite{schrijver}. 

Matematicamente:

\begin{equation}
    |MVC| + |MIS| = |V|
\end{equation}

Pelo Teorema de König, detalhado anteriormente, sabemos que em grafos bipartidos o tamanho da Cobertura Mínima de Vértices é exatamente igual ao tamanho do Emparelhamento Máximo ($|MVC| = |M|$). Consequentemente, o tamanho do Conjunto Independente Máximo pode ser deduzido diretamente do emparelhamento: $|MIS| = |V| - |M|$ \cite{jungnickel}.

Além da equivalência de cardinalidade, o Teorema de Gallai também fornece a dualidade estrutural: o Conjunto Independente Máximo é exatamente o complemento da Cobertura Mínima de Vértices. Portanto, para construir algoritmicamente o Conjunto Independente Máximo em um grafo bipartido, basta executar o seguinte procedimento:

\begin{enumerate}
    \item Calcula-se o Emparelhamento Máximo utilizando algoritmos como o de Hopcroft-Karp.\
    \item Extrai-se a Cobertura Mínima de Vértices $C$ através do método de busca por caminhos alternantes a partir dos vértices livres, definido na seção \ref{p:teorema_konig}.\
    \item O Conjunto Independente Máximo $I$ será composto por todos os vértices do grafo que não pertencem à cobertura. Ou seja, $I = V \setminus C$.
\end{enumerate}

\subsubsection{Algoritmos exatos}
\label{sec:algos_exatos_emp_card_max}

Por ser em um ambiente mais restrito, o problema de emparelhamento de cardinalidade máxima tem diversas soluções em tempo polinomial. A seguir listadas algumas delas:

\paragraph{Algoritmo do caminho aumentante}

O algoritmo de caminho aumentante se baseia em um lema proposto por Claude Berge, que afirma que um grafo G é máximo se e somente se não existe nenhum caminho aumentante em G. Este algoritmo é uma implementação direta deste lema, onde se busca repetidamente por caminhos aumentantes e se aumenta o emparelhamento até que nenhum caminho aumentante possa ser encontrado \cite{lawler}.

Começamos com um emparelhamento vazio $M=\emptyset$. Escolhemos um vértice livre $u$ no lado esquerdo do grafo bipartido e tentamos encontrar um caminho aumentante $P$ que começa em $u$ e termina em um vértice livre no lado direito do grafo, respeitando a regra da alternância entre arestas em $M$ e arestas fora de $M$. Se um caminho aumentante for encontrado, atualizamos o emparelhamento $M$ invertendo as arestas ao longo do caminho $P$. Repetimos esse processo até que não seja possível encontrar mais caminhos aumentantes. O emparelhamento resultante será o emparelhamento de cardinalidade máxima \cite{halim}.

Em seguida o algoritmo em pseudo-código:

\begin{algorithm}[H]
    \DontPrintSemicolon
    \SetAlgoLined

    % Variáveis globais
    \textbf{Global:} $match$ (vetor de emparelhamento), $vis$ (vetor de visitados), $AdjList$\;
    \bigskip

    % Procedimento Augment (DFS)
    \Fn{Aug($l$)}{
        \If{$vis[l] = 1$}{
            \KwRet 0\;
        }
        $vis[l] \leftarrow 1$\;
        
        \For{each $r \in AdjList[l]$}{
            % A lógica do caminho aumentante
            \If{$match[r] = -1$ \KwOr Aug($match[r]$)}{
                $match[r] \leftarrow l$\;
                \KwRet 1 \textbf{[comment: found 1 matching]}\;
            }
        }
        \KwRet 0 \textbf{[comment: no matching]}\;
    }

    \bigskip

    % Procedimento Principal
    \Fn{MCBM\_Algorithm($n$)}{
        $MCBM \leftarrow 0$\;
        $match \leftarrow -1$ (initialize for all $V$)\;
        
        \For{$l \leftarrow 0$ \KwTo $n-1$}{
            $vis \leftarrow 0$ (reset before each recursion)\;
            $MCBM \leftarrow MCBM + \text{Aug}(l)$\;
        }
        
        \KwRet $MCBM$\;
    }
    
    \caption{Algoritmo de caminhos aumentantes (MCBM) \cite{halim}}
    \label{alg:mcbm_augmenting}
\end{algorithm}

Como o algoritmo repete a busca por caminhos aumentantes para cada vértice do lado esquerdo do grafo, o tempo de execução total do algoritmo é $O(VE)$, onde $V$ é o número de vértices e $E$ é o número de arestas no grafo bipartido \cite{halim}.

\paragraph{Redução ao problema de fluxo máximo}

O problema de emparelhamento de cardinalidade máxima em grafos bipartidos pode ser eficientemente resolvido através de uma redução ao problema de fluxo máximo em redes. A ideia central é construir uma rede de fluxo a partir do grafo bipartido original, onde cada aresta do grafo bipartido é convertida em uma aresta com capacidade unitária na rede de fluxo. \cite{lawler}

Primeiro, adicionamos um vértice fonte $s$ e um vértice sumidouro $t$ à rede. Conectamos o vértice fonte $s$ a todos os vértices do conjunto esquerdo do grafo bipartido com arestas de capacidade 1. Em seguida, conectamos todos os vértices do conjunto direito do grafo bipartido ao vértice sumidouro $t$, também com arestas de capacidade 1. As arestas entre os conjuntos esquerdo e direito do grafo bipartido são mantidas na rede de fluxo, cada uma com capacidade 1 \cite{halim}.

Com isso, temos agora um grafo de fluxo onde o objetivo é encontrar o fluxo máximo do vértice fonte $s$ para o vértice sumidouro $t$. O valor do fluxo máximo encontrado nesta rede corresponde ao tamanho do emparelhamento máximo no grafo bipartido original. \cite{lawler}. Para encontrar o fluxo máximo, podemos utilizar algoritmos clássicos como o de Edmonds-Karp.

\paragraph{Algoritmo de Edmonds-Karp}
\label{p:edmonds-karp}

O algoritmo de \textbf{Edmonds-Karp} é o principal método para resolver esse tipo de problema. Ele é uma especialização do método de Ford-Fulkerson descrito na seção \ref{p:ford_fulkerson} que utiliza uma busca em largura (BFS) para encontrar sempre o caminho aumentante mais curto em número de arestas. Tem complexidade de tempo $O(VE^2)$:

\begin{algorithm}[H]
    \DontPrintSemicolon
    \SetAlgoLined

    \textbf{Global:} $res$ (matriz de capacidade residual), $mf, f, s, t$, $p$ (vetor de pais da BFS)\;
    \bigskip

    \Fn{BFS($s, t$)}{
        $dist \leftarrow \infty; \quad dist[s] \leftarrow 0$\;
        $Q \leftarrow \{s\}; \quad p \leftarrow \text{nil}$ (init vector with -1)\;
        \While{$Q \neq \emptyset$}{
            $u \leftarrow Q.\text{front}()$; $Q.\text{pop}()$\;
            \If{$u = t$}{
                \KwBreak (stop BFS)\;
            }
            \For{each $v \in V$}{
                \If{$res_{u, v} > 0$ \KwAnd $dist[v] = \infty$}{
                    $dist[v] \leftarrow dist[u] + 1$\;
                    $Q.\text{push}(v)$\;
                    $p[v] \leftarrow u$\;
                }
            }
        }
    }

    \bigskip

    \Fn{Augment($v, \text{minEdge}$)}{
        \eIf{$v = s$}{
            $f \leftarrow \text{minEdge}$\;
            \KwRet\;
        }{
            \If{$p[v] \neq -1$}{
                Augment($p[v], \min(\text{minEdge}, res_{p[v], v})$)\;
                $res_{p[v], v} \leftarrow res_{p[v], v} - f$\;
                $res_{v, p[v]} \leftarrow res_{v, p[v]} + f$\;
            }
        }
    }

    \bigskip

    \Fn{EdmondsKarp($s, t$)}{
        $mf \leftarrow 0$\;
        \While{\KwTrue}{
            $f \leftarrow 0$\;
            BFS($s, t$); \AlgComment{fills the parent vector $p$}\;
            Augment($t, \infty$); \AlgComment{find min edge weight $f$ in path}\;
            \If{$f = 0$}{
                \KwBreak (terminate)\;
            }
            $mf \leftarrow mf + f$\;
        }
        \KwRet $mf$\;
    }

    \caption{Algoritmo de Edmonds-Karp \cite{halim}}
    \label{alg:edmonds_karp}
\end{algorithm}

O pseudocódigo divide-se em três procedimentos:

\begin{itemize}
    \item \textbf{BFS:} Busca em largura a partir de $s$ na rede residual, percorrendo apenas arestas com capacidade residual positiva. Ao alcançar cada vértice pela primeira vez, registra seu pai em $p$; como a busca é em largura, o caminho registrado até $t$ é o de menor número de arestas.
    \item \textbf{Augment:} Percorre o caminho encontrado de $t$ até $s$, seguindo o vetor de pais. Na volta da recursão, $f$ guarda a menor capacidade residual do caminho (o gargalo), que é subtraída das arestas no sentido do caminho e somada às arestas reversas.
    \item \textbf{EdmondsKarp:} Repete a busca e o aumento enquanto existir caminho de $s$ a $t$. Quando a BFS não alcança $t$, nenhum fluxo é enviado ($f = 0$) e o fluxo acumulado $mf$ é máximo.
\end{itemize}

\paragraph{Algoritmo de Dinic}

Outra opção para resolver o problema de fluxo máximo é o algoritmo de Dinic, que é mais eficiente em muitos casos práticos, especialmente em grafos densos. O algoritmo de Dinic utiliza uma combinação de buscas em largura (BFS) para construir níveis na rede residual e buscas em profundidade (DFS) para encontrar caminhos aumentantes dentro desses níveis. O tempo de execução do algoritmo de Dinic é $O(V^2 E)$ em redes gerais e cai para $O(E \sqrt{V})$ em redes de capacidade unitária, como a obtida na redução do emparelhamento bipartido, o que o torna uma escolha preferida para muitos problemas de fluxo máximo \cite{halim}.

Este trabalho não tem uma implementação do algoritmo de Dinic. Ao invés disso, será apresentado o algoritmo de Hopcroft-Karp, que é uma variação especializada do algoritmo de Dinic para o problema de emparelhamento em grafos bipartidos, oferecendo uma solução mais eficiente para este caso específico.

\paragraph{Algoritmo de Hopcroft-Karp}

O algoritmo de Hopcroft-Karp é um algoritmo eficiente para encontrar o emparelhamento máximo em grafos bipartidos. Ele melhora o desempenho do algoritmo de caminhos aumentantes ao encontrar múltiplos caminhos aumentantes em cada iteração, em vez de apenas um. O tempo de execução do algoritmo de Hopcroft-Karp é $O(E \sqrt{V})$, tornando-o significativamente mais rápido do que o algoritmo de caminhos aumentantes simples, especialmente em grafos grandes \cite{halim}.

O algoritmo consiste em duas fases principais: a fase de construção de níveis e a fase de busca de caminhos aumentantes. Na fase de construção de níveis, uma busca em largura (BFS) é realizada a partir dos vértices livres no lado esquerdo do grafo bipartido para construir uma camada de níveis que ajuda a identificar os caminhos aumentantes mais curtos. Na fase de busca de caminhos aumentantes, uma busca em profundidade (DFS) é realizada para encontrar todos os caminhos aumentantes possíveis dentro da camada de níveis construída na fase anterior. Esses caminhos são então usados para aumentar o emparelhamento \cite{halim}.

A seguir, o pseudo-código do algoritmo de Hopcroft-Karp:

\begin{algorithm}[H]
    \DontPrintSemicolon
    \SetAlgoLined

    \caption{Algoritmo de Hopcroft-Karp \cite{burkard}}
    \label{alg:hopcroft_karp}

    % MAIN ALGORITHM
    \KwLet $G = (U, V; E)$ be a bipartite graph with initial (possibly empty) matching $M$;\;
    \KwLet $U_0$ contain all unmatched vertices of $U$;\label{hk:u0}\;
    \KwLet $V_0$ contain all unmatched vertices of $V$;\label{hk:v0}\;
    
    \While{$U_0 \neq \emptyset$\label{hk:phase}}{
        Construct\_layered\_graph;\label{hk:callbuild}\;
        Find\_set\_$\Delta M$;\label{hk:callfind}\;
        $M := M \ominus \Delta M$;\label{hk:symdiff}\;
        update sets $U_0, V_0$ of unmatched vertices\label{hk:upd}\;
    }

    \bigskip

    % PROCEDURE 1: CONSTRUCT LAYERED GRAPH
    \Fn{Construct\_layered\_graph}{
        $L_0 := U_0, k^* := k := 0$;\label{hk:l0}\;
        \While{$L_k \neq \emptyset$}{
            \For{each $i \in L_k$}{
                $N(i) := \{ j : [i, j] \in E \setminus M, j \notin L_1 \cup L_3 \cup \dots \cup L_{k-1} \}$;\label{hk:ni}
            }
            $L_{k+1} := \bigcup_{i \in L_k} N(i)$;\;
            
            \If{$L_{k+1} = \emptyset$\label{hk:stop}}{
                \KwStop \AlgComment{$M$ is a maximum matching}
            }
            
            \eIf{$L_{k+1} \cap V_0 \neq \emptyset$\label{hk:found}}{
                $k^* := k + 1$;\;
                $L_{k+2} := \emptyset$\;
            }{
                $L_{k+2} := \{ \overline{i} : [\overline{i}, j] \in M, j \in L_{k+1} \}$;\label{hk:lk2}
            }
            $k := k + 2$\;
        }
    }

    \bigskip

    % PROCEDURE 2: FIND SET DELTA M
    \Fn{Find\_set\_$\Delta M$}{
        \AlgComment{find a maximal set $\Delta M$ of disjoint augmenting paths}\;
        delete all vertices in $L_{k^*} \setminus V_0$;\label{hk:del}\;
        mark all remaining vertices as unscanned;\;
        $k := 1$; \AlgComment{path counter}\;
        
        \While{$L_0 \neq \emptyset$}{
            choose $x_0 := i \in L_0$ and delete it in $L_0$;\label{hk:x0}\;
            $\ell := 0$;\;
            
            \While{$\ell \ge 0$}{
                \While{$x_\ell$ has an unscanned neighbor in $L_{\ell+1}$\label{hk:adv}}{
                    choose unscanned neighbor $x_{\ell+1}$;\;
                    mark $x_{\ell+1}$ as scanned;\;
                    $\ell := \ell + 1$;\;
                    
                    \If{$\ell = k^*$\label{hk:path}}{
                        $P_k := (x_0, x_1, \dots, x_{k^*})$;\;
                        $k := k + 1$\;
                    }
                }
                \eIf{$\ell < k^*$\label{hk:back}}{
                    $\ell := \ell - 1$\;
                }{
                    $\ell := -1$\;
                }
            }
        }
        \KwRet $\Delta M := (P_1, P_2, \dots, P_{k-1})$\;
    }

\end{algorithm}

O pseudocódigo é organizado em um laço principal e dois procedimentos, explicados a seguir com referência às linhas do Algoritmo~\ref{alg:hopcroft_karp}:

\begin{itemize}
    \item \textbf{Laço principal (linhas~\ref{hk:u0}--\ref{hk:upd}):} $U_0$ e $V_0$ guardam os vértices livres de cada lado (linhas~\ref{hk:u0} e~\ref{hk:v0}). Cada iteração do laço da linha~\ref{hk:phase} é uma \textit{fase}: constrói-se o grafo em camadas (linha~\ref{hk:callbuild}), encontra-se um conjunto maximal $\Delta M$ de caminhos aumentantes mínimos e disjuntos (linha~\ref{hk:callfind}) e o emparelhamento é atualizado pela diferença simétrica $M \ominus \Delta M$ (linha~\ref{hk:symdiff}). Ao longo de cada caminho, as arestas emparelhadas e as não emparelhadas trocam de papel, de modo que $|M|$ aumenta em $|\Delta M|$.
    \item \textbf{Construção do grafo em camadas:} É uma busca em largura que parte de todos os vértices livres de $U$ ao mesmo tempo, que formam a camada $L_0$ (linha~\ref{hk:l0}). As camadas de índice ímpar contêm vértices de $V$, alcançados por arestas \textit{não} emparelhadas (linha~\ref{hk:ni}); a restrição $j \notin L_1 \cup \dots \cup L_{k-1}$ impede que um vértice seja revisitado. Se nenhum vértice novo é alcançado (linha~\ref{hk:stop}), não há caminho aumentante e, pelo lema de Berge, $M$ é máximo. Se a nova camada contém um vértice livre de $V$ (linha~\ref{hk:found}), foi encontrado o comprimento $k^*$ dos caminhos aumentantes mais curtos, e a construção termina. Caso contrário, a camada seguinte é formada pelos parceiros, em $M$, dos vértices recém-alcançados (linha~\ref{hk:lk2}): o caminho segue obrigatoriamente pela aresta emparelhada, o que garante a alternância.
    \item \textbf{Busca dos caminhos disjuntos:} Primeiro, descartam-se os vértices da última camada que não são livres (linha~\ref{hk:del}), pois nenhum caminho aumentante termina neles. Em seguida, uma busca em profundidade parte de cada vértice livre $x_0 \in L_0$ (linha~\ref{hk:x0}) e avança sempre para um vizinho ainda não visitado da camada seguinte (linha~\ref{hk:adv}). Ao atingir a camada $k^*$, o caminho $P_k = (x_0, \dots, x_{k^*})$ é registrado (linha~\ref{hk:path}) e a busca a partir dessa raiz termina; se a busca fica sem vizinhos antes disso, recua uma camada (linha~\ref{hk:back}). Como cada vértice é visitado no máximo uma vez por fase, os caminhos encontrados são disjuntos em vértices, e o conjunto $\Delta M$ é maximal.
    \item \textbf{Complexidade:} Cada fase examina cada aresta um número constante de vezes, com custo $O(E)$. Após $\sqrt{V}$ fases, todo caminho aumentante restante tem mais de $\sqrt{V}$ arestas; como os caminhos aumentantes de $M \ominus M^*$ são disjuntos em vértices, restam menos de $\sqrt{V}$ deles, e cada fase seguinte elimina ao menos um. O número de fases é, portanto, $O(\sqrt{V})$, o que leva à complexidade total $O(E\sqrt{V})$ \cite{burkard}.
\end{itemize}

\subsubsection{Heurísticas}
\label{sec:h_emp_card_max_bip}

Por mais que o algoritmo de Hopcroft-Karp descrito em \ref{alg:hopcroft_karp} seja considerado relativamente eficiente, com complexidade $O(E \sqrt{V})$, ele ainda não é tão viável para instâncias muito grandes do problema de emparelhamento de cardinalidade máxima. Por isso, foram propostos alguns métodos aproximados (heurísticas) com o objetivo de resolver o problema com complexidade linear $O(E)$.

\paragraph{Heurística gulosa simples}
\label{p:h_gulosa_simples}

Um método que itera sobre as arestas do grafo em ordem arbitrária. Se a aresta $(u,v)$ conecta dois vértices livres, adicione essa aresta ao emparelhamento e remova os vértices $u$ e $v$. Tem complexidade de tempo linear $O(E)$ e garante um emparelhamento maximal, que conforme visto na seção \ref{p:emp_max_maximal}, é um emparelhamento ao qual não é possível adicionar mais arestas, mas que não é necessariamente o de maior cardinalidade. A solução também garante um emparelhamento de pelo menos $1/2$ (metade) da cardinalidade máxima \cite{vazirani2001}.

\begin{algorithm}[H]
    \DontPrintSemicolon
    \SetAlgoLined

    \caption{Heurística gulosa simples \cite{vazirani2001}}
    \label{alg:greedy_mcm}

    \KwIn{bipartite graph $G = (U, V; E)$}
    \KwOut{maximal matching $M$}
    $M := \emptyset$\;
    \For{each edge $(u, v) \in E$, in arbitrary order}{
        \If{$u$ and $v$ are both unmatched}{
            $M := M \cup \{(u, v)\}$\;
            mark $u$ and $v$ as matched\;
        }
    }
    \KwRet $M$\;
\end{algorithm}

\paragraph{Heurística de Karp-Sipser}

Heurística descrita por \citeonline{karp-sipser} para grafos esparsos que toma decisões com base na topologia do grafo ao invés de decisões aleatórias, tornando essas decisões mais seguras.

O algoritmo funciona reduzindo dinamicamente o grafo $G=(V,E)$ através da aplicação repetida de duas regras, nesta ordem de prioridade:

\begin{enumerate}
    \item \textbf{Redução de folhas de grau 1:} Se existir algum vértice $v$ no grafo atual com \textbf{grau 1} (apenas um vizinho $u$), a aresta $(u, v)$ é selecionada imediatamente para o emparelhamento. Isso é seguro pois, se $u$ é o único vizinho de $v$, ele também é a única opção para emparelhar $v$. Se isso não for feito, $v$ ficará isolado.
    \item \textbf{Escolha Aleatória:} Se \textbf{não} existirem vértices de grau 1 (ou seja, todos os vértices têm grau $\ge 2$), o algoritmo escolhe uma aresta $(u, v)$ aleatoriamente, adiciona ao emparelhamento e remove $u$ e $v$ do grafo. 
    
    Em grafos esparsos, isso tende a quebrar ciclos e formar novas folhas (vértices de grau 1), permitindo que o algoritmo volte a usar a regra 1. Porém, em grafos densos, essa escolha aleatória pode comprometer a precisão da solução.
\end{enumerate}

Para grafos esparsos, a regra 2 é raramente aplicada, pois a remoção de vértices na regra 1 frequentemente cria novas folhas, provocando um efeito cascata, fazendo com que a maior parte do grafo seja resolvida de maneira ótima. Com estruturas de dados adequadas, a complexidade do algoritmo é linear ($O(E)$). A implementação avaliada neste trabalho, por simplicidade, localiza os vértices de grau 1 por varredura, o que eleva sua complexidade a $O(V \cdot E)$.

\begin{algorithm}[H]
    \DontPrintSemicolon
    \SetAlgoLined

    \caption{Heurística de Karp-Sipser \cite{karp-sipser}}
    \label{alg:karp_sipser}

    \KwIn{bipartite graph $G = (U, V; E)$}
    \KwOut{maximal matching $M$}
    $M := \emptyset$\;
    \While{$E \neq \emptyset$}{
        \eIf{there is a vertex $v$ with $\deg(v) = 1$}{
            let $u$ be the only neighbor of $v$\;
            \AlgComment{safe: some maximum matching contains $(u, v)$}\;
        }{
            choose an arbitrary edge $(u, v) \in E$\;
            \AlgComment{the only step that may lose optimality}\;
        }
        $M := M \cup \{(u, v)\}$\;
        remove $u$, $v$ and all their incident edges from $G$\;
    }
    \KwRet $M$\;
\end{algorithm}

\begin{table}[htbp]
    \centering
    \caption{Algoritmos para emparelhamento de cardinalidade máxima em grafos bipartidos}
    \label{tab:emp_card_max_bip_algos}
    \renewcommand{\arraystretch}{1.2} % Adds a little vertical padding
    
    % Structure: l = left align, c = center, X = auto-wrap paragraph
    \begin{tabularx}{\textwidth}{@{} l l l l X @{}}
        \toprule
        \textbf{Algoritmo} & \textbf{Complexidade} & \textbf{Tipo} & \textbf{Avaliado} & \textbf{Descrição} \\
        \midrule
        
        Caminhos aumentantes & $O(VE)$ & Exato & Sim & 
        Busca iterativa (DFS/BFS) por caminhos alternantes entre vértices livres, invertendo arestas para aumentar a cardinalidade. \\
        \addlinespace 
        
        Edmonds-Karp & $O(VE^2)$ & Exato & Sim & 
        Implementação de Ford-Fulkerson que usa BFS para encontrar sempre o caminho aumentante mais curto (em nº de arestas) na rede residual. \\
        \addlinespace
        
        Dinic & $O(V^2E)$ & Exato & Não & 
        Algoritmo de fluxo que constrói grafos de níveis (BFS) e satura múltiplos caminhos simultaneamente via fluxo de bloqueio (DFS). Em redes de capacidade unitária, como a do emparelhamento bipartido, cai para $O(E\sqrt{V})$. \\
        \addlinespace
        
        Hopcroft-Karp & $O(E\sqrt{V})$ & Exato & Sim & 
        Busca múltiplos caminhos aumentantes disjuntos por fase usando BFS e DFS. \\
        \addlinespace
        
        Gulosa simples & $O(E)$ & Heurística & Sim & 
        Itera sobre arestas arbitrariamente, selecionando aquelas que conectam dois vértices livres. \\
        \addlinespace

        Karp-Sipser & $O(E)$ & Heurística & Sim* & 
        Prioriza remoção de vértices de grau 1; usa escolha aleatória quando não há folhas. \\
        
        \bottomrule
    \end{tabularx}
    \par\smallskip
    {\raggedright\footnotesize \textit{Avaliado}: incluído no ambiente de testes; a origem de cada implementação consta na Tabela~\ref{tab:proveniencia}. * A implementação avaliada localiza os vértices de grau 1 por varredura simples, com complexidade $O(V \cdot E)$.\par}
\end{table}

% TODO: possivelmente definir aplicações? exemplo: min-path-cover.md
% ----------------------------------------------------------------------------------------------
% 3.2.1. Problema de atribuição (Assignment Problem)
% ----------------------------------------------------------------------------------------------
\newpage
\subsection{Problema de atribuição (Assignment Problem)}
\label{sec:assignment_problem}

\subsubsection{Descrição do problema}

O problema de atribuição, também conhecido como problema do emparelhamento ponderado, consiste em encontrar uma combinação ótima de atribuições entre dois conjuntos disjuntos, minimizando o custo total associado a essas atribuições. Exemplo: Considere $N$ trabalhadores e $N$ tarefas, onde cada trabalhador pode ser designado a exatamente uma tarefa, e cada tarefa deve ser atribuída a exatamente um trabalhador. O custo de atribuir o trabalhador i à tarefa j é representado por uma matriz de custos $W = [w_{ij}]$ \cite{lawler}. O objetivo é encontrar um conjunto de atribuições que minimize o custo total.

\subsubsection{Propriedades}

\subsubsection{Algoritmos exatos}

% \paragraph{Método Húngaro}
\paragraph{Método Húngaro}

O método Húngaro é um algoritmo eficiente para resolver o problema de atribuição em tempo polinomial. Ele foi desenvolvido por Harold Kuhn em 1955 e é baseado no trabalho de Dénes Kőnig e Jenő Egerváry sobre emparelhamentos em grafos bipartidos \cite{jungnickel}.

O método se baseia em manter um conjunto de potenciais para os vértices dos dois conjuntos do grafo bipartido, sendo $u_i$ o potencial do vértice i no conjunto esquerdo e $v_j$ o potencial do vértice j no conjunto direito. A regra que o algoritmo mantém é que para cada aresta $(i,j)$, a soma dos potenciais deve ser maior ou igual ao custo da aresta, ou seja, $u_i + v_j \geq w_{ij}$. 

Cria-se um subgrafo de igualdade $H_{u,v}$ contendo apenas as arestas onde a soma dos potenciais é igual ao custo da aresta, ou seja, $u_i + v_j = w_{ij}$. O algoritmo então tenta encontrar um emparelhamento máximo neste subgrafo de igualdade. Se o emparelhamento encontrado cobre todos os vértices do conjunto esquerdo, então ele é ótimo para o problema de atribuição original. Caso contrário, o algoritmo ajusta os potenciais para criar novas arestas de igualdade e repete o processo até encontrar um emparelhamento máximo que cubra todos os vértices do conjunto esquerdo \cite{jungnickel}.

Em seguida o pseudo-código do método Húngaro:

\begin{algorithm}[H]
    \DontPrintSemicolon
    \SetAlgoLined

    \Fn{HUNGARIAN($n, w; mate$)}{
        
        % Lines 1-3
        \lFor{$v \in V$}{ $mate(v) \leftarrow 0$ \label{hu:init} }
        \lFor{$i = 1$ \KwTo $n$}{ $u_i \leftarrow \max \{ w_{ij} : j=1,\dots,n \}; v_i \leftarrow 0$ }
        $nrex \leftarrow n$\label{hu:nrex}\;

        % Line 4: While Loop
        \While{$nrex \neq 0$\label{hu:phase}}{
            
            % Line 5
            \lFor{$i = 1$ \KwTo $n$}{ $m(i) \leftarrow \text{false}; p(i) \leftarrow 0; \delta_i \leftarrow \infty$ \label{hu:reset} }
            $aug \leftarrow \text{false}; Q \leftarrow \{ i \in S : mate(i) = 0 \}$\label{hu:roots}\;
            
            % Line 7: Repeat Loop
            \Repeat{$aug = \text{true}$}{
                remove an arbitrary vertex $i$ from $Q$; $m(i) \leftarrow \text{true}; j \leftarrow 1$\label{hu:scan}\;
                
                % Line 9: Inner While
                \While{$aug = \text{false}$ \KwAnd $j \leq n$}{
                    \If{$mate(i) \neq j'$\label{hu:skip}}{
                         \If{$u_i + v_j - w_{ij} < \delta_j$\label{hu:slack}}{
                            $\delta_j \leftarrow u_i + v_j - w_{ij}; p(j) \leftarrow i$\;
                            \If{$\delta_j = 0$\label{hu:tight}}{
                                \eIf{$mate(j') = 0$}{
                                    AUGMENT($mate, p, j'; mate$)\label{hu:aug1}\;
                                    $aug \leftarrow \text{true}; nrex \leftarrow nrex - 1$\;
                                }{
                                    $Q \leftarrow Q \cup \{ mate(j') \}$\label{hu:grow}\;
                                }
                            }
                        }
                    }
                    $j \leftarrow j + 1$\;
                }
                
                % Line 24: Check if augmentation failed and Q is empty
                \If{$aug = \text{false}$ \KwAnd $Q = \emptyset$\label{hu:dual}}{
                    $J \leftarrow \{ i \in S : m(i) = \text{true} \}; K \leftarrow \{ j' \in T : \delta_j = 0 \}$\;
                    $\delta \leftarrow \min \{ \delta_j : j' \in T \setminus K \}$\label{hu:delta}\;
                    \lFor{$i \in J$}{ $u_i \leftarrow u_i - \delta$ }
                    \lFor{$j' \in K$}{ $v_j \leftarrow v_j + \delta$ }
                    \lFor{$j' \in T \setminus K$}{ $\delta_j \leftarrow \delta_j - \delta$ \label{hu:dj} }
                    $X \leftarrow \{ j' \in T \setminus K : \delta_j = 0 \}$\label{hu:X}\;
                    
                    \eIf{$mate(j') \neq 0$ for all $j' \in X$}{
                        \lFor{$j' \in X$}{ $Q \leftarrow Q \cup \{ mate(j') \}$ }
                    }{
                        choose $j' \in X$ with $mate(j') = 0$\label{hu:aug2}\;
                        AUGMENT($mate, p, j'; mate$)\;
                        $aug \leftarrow \text{true}; nrex \leftarrow nrex - 1$\;
                    }
                }
            }
        }
    }
    \caption{Algoritmo Húngaro \cite{jungnickel}}
    \label{alg:hungarian}
\end{algorithm}

\begin{algorithm}[H]
    \DontPrintSemicolon
    \SetAlgoLined

    \Fn{AUGMENT($mate, p, j'; mate$)}{
        \Repeat{$next = 0$}{
            $i \leftarrow p(j); mate(j') \leftarrow i; next \leftarrow mate(i); mate(i) \leftarrow j'$\;
            \If{$next \neq 0$}{
                $j' \leftarrow next$\;
            }
        }
    }
    \caption{Método AUGMENT \cite{jungnickel}}
    \label{alg:augment}
\end{algorithm}

Na versão de \citeonline{jungnickel}, o algoritmo \textit{maximiza} o peso do emparelhamento; para minimizar custos, basta usar $w_{ij} = -c_{ij}$. As linhas da matriz formam o conjunto $S = \{1, \dots, n\}$ e as colunas, o conjunto $T = \{1', \dots, n'\}$, em que $j'$ denota a coluna $j$ como vértice do grafo. A seguir, cada parte do Algoritmo~\ref{alg:hungarian} é explicada com referência às suas linhas:

\begin{itemize}
    \item \textbf{Inicialização (linhas~\ref{hu:init}--\ref{hu:nrex}):} Nenhum vértice começa emparelhado. Cada potencial de linha $u_i$ recebe o maior peso da linha $i$, e os potenciais de coluna $v_j$ recebem zero, o que satisfaz $u_i + v_j \geq w_{ij}$ para toda aresta. O contador $nrex$ guarda o número de linhas ainda livres.
    \item \textbf{Fases (linha~\ref{hu:phase}):} Cada iteração do laço externo aumenta o emparelhamento em uma aresta, até que todas as linhas estejam emparelhadas. No início da fase, as marcas são reiniciadas (linha~\ref{hu:reset}): $m(i)$ indica se a linha $i$ já entrou na árvore de busca, $p(j)$ guarda a linha pela qual a coluna $j$ foi alcançada e $\delta_j$ guarda a menor folga $u_i + v_j - w_{ij}$ da coluna $j$ em relação às linhas da árvore. A fila $Q$ começa com as linhas livres, que são as raízes da busca (linha~\ref{hu:roots}).
    \item \textbf{Busca no subgrafo de igualdade:} Uma linha $i$ é retirada da fila e entra na árvore (linha~\ref{hu:scan}). Para cada coluna que não é a sua parceira atual (linha~\ref{hu:skip}), calcula-se a folga da aresta $(i, j)$ e, se ela for menor que $\delta_j$, a coluna passa a ter $i$ como predecessora (linha~\ref{hu:slack}). Folga nula significa que a aresta pertence ao subgrafo de igualdade $H_{u,v}$ (linha~\ref{hu:tight}). Nesse caso, se a coluna está livre, foi encontrado um caminho aumentante, e o procedimento AUGMENT o aplica (linha~\ref{hu:aug1}); senão, a parceira da coluna entra na fila (linha~\ref{hu:grow}), e a árvore cresce pela aresta emparelhada.
    \item \textbf{Ajuste dos potenciais (linhas~\ref{hu:dual}--\ref{hu:aug2}):} Se a fila esvazia sem que um caminho aumentante seja encontrado, as linhas e colunas da árvore formam os conjuntos $J$ e $K$. Calcula-se $\delta$, a menor folga entre as colunas fora da árvore (linha~\ref{hu:delta}). Diminuir $u_i$ em $\delta$ nas linhas da árvore e aumentar $v_j$ em $\delta$ nas colunas da árvore preserva as arestas da árvore no subgrafo de igualdade e mantém a viabilidade dos potenciais. Ao mesmo tempo, a folga das colunas de fora diminui em $\delta$ (linha~\ref{hu:dj}), o que coloca ao menos uma nova aresta no subgrafo de igualdade (o conjunto $X$, linha~\ref{hu:X}). Se todas as colunas de $X$ estão emparelhadas, suas parceiras entram na fila e a busca continua; senão, uma coluna livre de $X$ completa um caminho aumentante (linha~\ref{hu:aug2}).
    \item \textbf{Aumento (Algoritmo~\ref{alg:augment}):} A partir da coluna livre encontrada, o procedimento AUGMENT segue os predecessores $p$ de volta até a raiz, invertendo o emparelhamento ao longo do caminho.
    \item \textbf{Complexidade:} Em cada fase, cada linha entra na árvore no máximo uma vez, e sua varredura custa $O(n)$. Cada ajuste de potenciais custa $O(n)$ e acrescenta ao menos uma coluna à árvore, de modo que há no máximo $n$ ajustes por fase. Cada fase custa, portanto, $O(n^2)$, e as $n$ fases resultam em $O(n^3)$.
\end{itemize}

\paragraph{Jonker-Volgenant} 

O algoritmo de Jonker-Volgenant é uma evolução otimizada para a resolução do problema de atribuição linear, proposto por Roy Jonker e Anton Volgenant em 1987. Ele foi desenvolvido para superar o desempenho prático do método Húngaro, especialmente em grafos densos, combinando a teoria de dualidade com estratégias de inicialização e busca mais eficientes \cite{jonker}.

Assim como o método Húngaro, este algoritmo baseia-se na manutenção de potenciais duais ($u_i$ e $v_j$) e respeita as condições de folga complementar. No entanto, sua principal distinção conceitual é o uso da estratégia de Caminho Aumentante Mais Curto (Shortest Augmenting Path). Ao invés de construir emparelhamentos máximos em fases distintas, o algoritmo foca em encontrar o caminho de custo mínimo que conecta uma linha não atribuída a uma coluna livre.

O processo inicia com heurísticas avançadas (conhecidas como redução de colunas e transferência de redução) para resolver rapidamente as atribuições triviais. Para os vértices restantes, o algoritmo realiza uma busca similar ao algoritmo de Dijkstra: ele explora o grafo para encontrar o caminho aumentante mais curto, atualizando os potenciais duais simultaneamente durante a busca. Isso garante que, a cada iteração, o emparelhamento seja aumentado pelo menor custo possível até que a solução ótima completa seja atingida.

A implementação completa deste algoritmo é extensa; o Algoritmo~\ref{alg:jv} apresenta suas quatro fases em alto nível. Sua eficiência prática o torna uma escolha preferida para muitos problemas de atribuição em aplicações reais, especialmente quando comparado ao método Húngaro tradicional \cite{jungnickel}.

\begin{algorithm}[H]
    \DontPrintSemicolon
    \SetAlgoLined

    \caption{Algoritmo de Jonker-Volgenant (LAPJV), em alto nível \cite{jonker}}
    \label{alg:jv}

    \KwIn{cost matrix $C = (c_{ij})_{n \times n}$}
    \KwOut{optimal assignment $x$ (row $i$ is assigned to column $x_i$) and column duals $v$}

    \AlgComment{phase 1: column reduction}\;
    \For{$j := n$ \KwTo $1$}{
        $i^* := \arg\min_i c_{ij}$; $v_j := c_{i^* j}$\;
        \lIf{row $i^*$ is unassigned}{assign $i^*$ to $j$}
    }

    \AlgComment{phase 2: reduction transfer}\;
    \For{each assigned row $i$ that is the minimum of a single column $j_1 = x_i$}{
        $v_{j_1} := v_{j_1} - \min_{j \neq j_1} (c_{ij} - v_j)$\;
    }

    \AlgComment{phase 3: augmenting row reduction (applied twice)}\;
    \For{each free row $i$}{
        let $j_1, j_2$ be the columns with the smallest and second smallest values $u_1 \le u_2$ of $c_{ij} - v_j$\;
        \eIf{$u_1 < u_2$}{
            $v_{j_1} := v_{j_1} - (u_2 - u_1)$\;
        }{
            \lIf{$j_1$ is assigned}{$j_1 := j_2$}
        }
        assign $i$ to $j_1$; the row previously assigned to $j_1$, if any, becomes free\;
    }

    \AlgComment{phase 4: augmentation by shortest paths}\;
    \For{each remaining free row $i$}{
        find a shortest augmenting path from $i$ to a free column with a Dijkstra-like search on the reduced costs $c_{kj} - v_j$\;
        $v_j := v_j + d_j - \mu$ for every scanned column $j$, where $d_j$ is its distance and $\mu$ the length of the path\;
        augment the assignment along the path\;
    }
    \KwRet $x$, $v$\;
\end{algorithm}

\paragraph{Redução para problema de fluxo de custo mínimo}
\label{p:reduction_to_minimal_flow}

Similarmente ao que foi feito para o problema de emparelhamento de cardinalidade máxima, o problema de atribuição pode ser resolvido através de uma redução ao problema de fluxo de custo mínimo. Assim como para o problema anterior, essa redução é feita criando dois novos vértices: um vértice fonte $s$ e um vértice sorvedouro $t$. Em seguida, são adicionadas arestas do vértice fonte $s$ para cada vértice de um dos conjuntos do grafo bipartido, com capacidade 1 e custo 0. Também são adicionadas arestas de cada vértice do outro conjunto para o vértice sorvedouro $t$, também com capacidade 1 e custo 0. Finalmente, as arestas do grafo bipartido original são adicionadas com capacidade 1 e custo igual ao custo de atribuição correspondente.

Para resolver problemas de fluxo podemos utilizar algoritmos como o de caminhos de custo mínimo sucessivos (Successive Shortest Paths) ou o algoritmo de cancelamento de ciclos (Cycle canceling) \cite{schrijver}.

% TODO: Incluir também Out-of-kilter, Simplex, e cost-scaling

O primeiro algoritmo, o de caminhos de custo mínimo sucessivos, foi inicialmente proposto por Jewell, Busacker e Gowen e é uma adaptação do algoritmo de Ford-Fulkerson para incorporar o custo das arestas. Em poucas palavras, ele encontra o caminho de $s$ para $t$ de custo mínimo, envia o máximo fluxo possível ao longo desse caminho, e repete esse processo até que não seja mais possível aumentar o fluxo sem violar as capacidades das arestas \cite{schrijver}.

\begin{algorithm}[H]
    \DontPrintSemicolon
    \SetAlgoLined

    \Fn{OPTFLOW($G, c, s, t, \gamma, v; f, \text{sol}$)}{
        % Line 1: Single-line for loop
        \lFor{$e \in E$}{ $f(e) \leftarrow 0$ }
        
        % Line 2
        $\text{sol} \leftarrow \text{true}, \text{val} \leftarrow 0$\;

        % Line 3: While loop
        \While{$\text{sol} = \text{true}$ \KwAnd $\text{val} < v$}{
            
            % Line 4: Descriptive text
            construct the auxiliary network $N' = (G', c', s, t)$ with cost function $\gamma'$\;

            % Line 5: If-Else structure
            \eIf{$t$ is not accessible from $s$ in $G'$}{
                % Line 6
                $\text{sol} \leftarrow \text{false}$\;
            }{
                % Line 7: Descriptive text
                determine a shortest path $P$ from $s$ to $t$ in $(G', \gamma')$\;
                
                % Line 8: Multiple assignments
                $\delta \leftarrow \min \{c'(e) : e \in P \}; \delta' \leftarrow \min(\delta, v - \text{val}); \text{val} \leftarrow \text{val} + \delta'$\;
                
                % Line 9
                augment $f$ along $P$ by $\delta'$\;
            }
        }
    }
    \caption{Algoritmo de caminhos de custo mínimo sucessivos \cite{jungnickel}}
    \label{alg:optflow}
\end{algorithm}

 A complexidade do algoritmo é dependente do método utilizado para encontrar o caminho $P$. A implementação padrão utilizaria o algoritmo de Dijkstra, que tem uma complexidade de $O(E + V \log V)$ em uma implementação com filas de prioridade, onde $V$ é o número de vértices e $E$ o número de arestas. Porém, como o grafo residual pode conter arestas com custos negativos, essa implementação não funcionaria corretamente. Para contornar esse problema, temos duas opções principais:

 \begin{itemize}
    \item \textbf{Utilizar Bellman-Ford:} A solução mais simples e direta é utilizar, ao invés do algoritmo de Dijkstra, o algoritmo de Bellman-Ford para encontrar o caminho de custo mínimo $P$. O Bellman-Ford é mais lento, com uma complexidade de $O(VE)$, mas lida corretamente com arestas de custo negativo.
    \item \textbf{Utilizar a técnica de Potenciais:} Permite que utilizemos o algoritmo de Dijkstra mesmo na presença de arestas de custo negativo.
    
    A cada iteração do loop, encontra-se o caminho de custo mínimo $P$ utilizando Dijkstra com os custos reduzidos. Após encontrar o caminho, os potenciais são atualizados para todos os vértices $v$ alcançáveis a partir de $s$ no grafo residual, ajustando-os conforme a distância mínima encontrada pelo Dijkstra, de acordo com a fórmula $h(v) \leftarrow h(v) + d(s,v)$, onde $d(s,v)$ é a distância mínima de $s$ até $v$ no grafo residual com custos reduzidos.

    Abaixo uma adaptação do algoritmo \ref{alg:optflow} utilizando a técnica de potenciais:

    \begin{algorithm}[H]
    \DontPrintSemicolon
    \SetAlgoLined

    \Fn{OPTFLOW\_DIJKSTRA($G, c, s, t, \gamma, v; f, \text{sol}$)}{
        % Inicialização do Fluxo
        \lFor{$e \in E$}{ $f(e) \leftarrow 0$ }
        
        % Inicialização dos Potenciais (assumindo custos iniciais não negativos)
        \lFor{$x \in V$}{ $\pi(x) \leftarrow 0$ }

        $\text{sol} \leftarrow \text{true}, \text{val} \leftarrow 0$\;

        \While{$\text{sol} = \text{true}$ \KwAnd $\text{val} < v$}{
            
            construct the auxiliary network $N' = (G', c', s, t)$\;
            
            % Definição do Custo Reduzido para o Dijkstra
            \For{$e=(u, v) \in E(G')$}{
                define reduced cost: $w(e) \leftarrow \gamma(e) + \pi(u) - \pi(v)$\;
            }

            % Execução do Dijkstra (Lógica omitida, foco no input/output)
            Run Dijkstra on $N'$ using weights $w$ to find shortest path distances $d(\cdot)$ from $s$\;

            \eIf{$t$ is not accessible from $s$ (i.e., $d(t) = \infty$)}{
                $\text{sol} \leftarrow \text{false}$\;
            }{
                % Atualização dos Potenciais (Passo Crucial)
                \lFor{$x \in V$ such that $d(x) < \infty$}{ $\pi(x) \leftarrow \pi(x) + d(x)$ }
                
                Identify shortest path $P$ from $s$ to $t$ based on $d(\cdot)$\;
                
                % O resto segue o algoritmo original, usando as capacidades residuais c'
                $\delta \leftarrow \min \{c'(e) : e \in P \}; \delta' \leftarrow \min(\delta, v - \text{val}); \text{val} \leftarrow \text{val} + \delta'$\;
                
                augment $f$ along $P$ by $\delta'$\;
            }
        }
    }
    \caption{Adaptação do algoritmo \ref{alg:optflow} utilizando Dijkstra com técnica de potenciais}
    \label{alg:optflow_dijkstra}
\end{algorithm}

A complexidade do algoritmo utilizando Dijkstra com a técnica de potenciais é $O(V (E + V \log V))$, onde $V$ é o número de vértices e $E$ o número de arestas do grafo. Isso ocorre porque, em cada iteração do loop principal, executamos o algoritmo de Dijkstra, que tem complexidade $O(E + V \log V)$, e o número máximo de iterações é limitado por $V$, o número de vértices no grafo bipartido.

\end{itemize}

\subsubsection{Heurísticas}

Assim como no problema de emparelhamento de cardinalidade máxima, as soluções exatas, apesar de apresentarem tempo polinomial, são inviáveis para instâncias muito grandes. A seguir algumas heurísticas que objetivam resolver o problema de atribuição em tempo até mesmo linear ($O(E)$).

\paragraph{Heurística gulosa ordenada}
\label{p:gulosa_ordenada}

Similar à heurística gulosa simples descrita na seção \ref{p:h_gulosa_simples}, com a diferença de que a ordem das arestas não é mais arbitrária, e sim decrescente por peso. Tem a mesma garantia de metade do peso ótimo, mas tem complexidade $O(E \log E)$ por conta da ordenação.

% É EXATO
% \paragraph{Algoritmo de leilão com $\epsilon$-scaling}

% O Algoritmo de Leilão (\textbf{Auction Algorithm}), descrito em \citeonline{berstekas} e \citeonline{berstekas1988}, resolve o problema da atribuição maximizando o benefício total através de um processo iterativo de lances e aumentos de preços, análogo a um leilão do mundo real.

% Considere $N$ pessoas e $N$ objetos.
% \begin{itemize}
%     \item $a_{ij}$: O valor (benefício) que a pessoa $i$ atribui ao objeto $j$.
%     \item $p_j$: O preço atual do objeto $j$.
% \end{itemize}

% O objetivo de cada pessoa $i$ é maximizar seu próprio \textbf{lucro líquido}:
% \begin{equation}
%     \text{Lucro}_{ij} = a_{ij} - p_j
% \end{equation}

% Em um equilíbrio de mercado (solução ótima), cada pessoa deve estar designada a um objeto que maximize seu lucro pessoal, e os preços devem ser tais que todos estejam satisfeitos. Isso corresponde, na teoria de dualidade, à condição de \textbf{Folga Complementar}.

% Se exigirmos que cada pessoa pegue *estritamente* o melhor objeto, o algoritmo pode entrar em um loop infinito de "guerra de preços" com incrementos nulos ou infinitesimais. Para garantir a convergência, Bertsekas introduziu o conceito de $\epsilon$-CS (Folga Complementar Epsilon).

% Uma atribuição é $\epsilon$-ótima se cada pessoa $i$ designada ao objeto $j$ estiver a uma distância $\epsilon$ de seu lucro máximo:

% \begin{equation}
%     (a_{ij} - p_j) \ge \max_{k} (a_{ik} - p_k) - \epsilon
% \end{equation}

% O Passo do Leilão:
% \begin{enumerate}
%     \item Uma pessoa não designada $i$ escolhe o objeto $j^*$ que oferece o maior lucro.
%     \item Ela calcula a diferença entre o lucro do melhor objeto e o lucro do segundo melhor objeto ($\gamma$).
%     \item Ela faz um lance para o objeto $j^*$, aumentando seu preço $p_{j^*}$ em:
%         $$\Delta p = \gamma + \epsilon$$
%     \item A pessoa $i$ ganha o objeto, e quem estava com ele (se houvesse alguém) torna-se não designado e volta à fila para dar lances.
% \end{enumerate}

% O termo $\epsilon$ é crucial: ele garante que o preço suba pelo menos um valor fixo a cada lance, impedindo loops infinitos e forçando o algoritmo a terminar. A escolha do valor de $\epsilon$ é muito importante pois define a eficiência do algoritmo:

% \begin{itemize}
%     \item Se $\epsilon$ for \textbf{grande}: O algoritmo termina muito rápido (poucos lances para resolver conflitos), mas a solução é de baixa qualidade (distante do ótimo).
%     \item Se $\epsilon$ for \textbf{pequeno}: A solução é muito próxima ou igual à ótima, mas o algoritmo é lento (muitos incrementos pequenos de preço necessários).
% \end{itemize}

% O conceito de $\epsilon$-scaling consiste em um refinamento sucessivo do valor de $\epsilon$:

% \begin{enumerate}
%     \item Comece com um $\epsilon$ grande (ex: metade do maior peso da aresta).
%     \item Rode o leilão até convergir (todos emparelhados).
%     \item Reduza o valor de $\epsilon$ (ex: $\epsilon \leftarrow \epsilon / 4$).
%     \item Mantenha os preços ($p_j$) e a atribuição da fase anterior. Reinicie o leilão com o novo $\epsilon$.
% \end{enumerate}

% Ao reutilizar os preços, as fases subsequentes convergem muito mais rápido do que se começassem do zero, pois os preços já estão "quase" certos. O processo repete-se até que $\epsilon < 1/N$, momento em que a solução obtida é comprovadamente a ótima inteira (para pesos inteiros).

% A complexidade do algoritmo (usando $\epsilon$-scaling) é de $O(N \cdot E \log(C))$, onde $C$ é o maior custo da aresta. Uma melhoria significativa dos algoritmos $O(N^3)$, como o Método Húngaro. Além disso, o algoritmo é paralelizável, pois permite que as pessoas façam lances simultaneamente, sendo ideal para aplicações onde é possível utilizar GPUs ou clusters distribuídos para o processamento.

% TODO: (pós TCC I): aproximação via Sinkhorn
% #### F. Aproximação via Sinkhorn (Soft Assignment)
% Muito popular atualmente em *Deep Learning* e Visão Computacional.
% * **Funcionamento:** Relaxa a restrição de "inteiro" (0 ou 1) para contínuo (probabilidade). Aplica normalização iterativa de linhas e colunas na matriz de custos (Algoritmo de Sinkhorn-Knopp).
% * **Resultado:** Gera uma matriz duplamente estocástica que aproxima a permutação ótima. Uma passagem final gulosa converte isso em uma atribuição discreta.
% * **Complexidade:** $O(k \cdot N^2)$ onde $k$ é pequeno. Muito rápido em hardware paralelo.

\newpage
\begin{table}[ht]
    \centering
    \caption{Algoritmos para problema de atribuição}
    \label{tab:atribuicao_algos}
    \renewcommand{\arraystretch}{1.2} % Adds a little vertical padding
    
    % Structure: l = left align, c = center, X = auto-wrap paragraph
    \begin{tabularx}{\textwidth}{@{} l l l X @{}}
        \toprule
        \textbf{Algoritmo} & \textbf{Complexidade de tempo} & \textbf{Tipo} & \textbf{Descrição} \\
        \midrule
        
        Método Húngaro & $O(n^3)$ & Exato & 
        Método primal-dual que mantém folga complementar para custo mínimo em bipartidos. \\
        \addlinespace 
        
        Jonker-Volgenant & $O(n^3)$ & Exato & 
        Algoritmo primal-dual eficiente que combina redução de custos inicial e busca de caminhos aumentantes com atualização otimizada de duais. \\
        \addlinespace
        
        Custo mínimo (Bellman-Ford) & $O(V^2E)$ & Exato & 
        Método de caminhos sucessivos que usa Bellman-Ford para encontrar caminhos na rede residual, suportando custos negativos iniciais. \\
        \addlinespace
        
        Custo mínimo (Dijkstra) & $O(V(E+V \log V))$ & Exato & 
        Utiliza potenciais de vértices para manter custos reduzidos não-negativos, permitindo o uso do eficiente Dijkstra para aumentos sucessivos. \\
        \addlinespace
        
        Gulosa ordenada & $O(E \log E)$ & Heurística & 
        Ordena as arestas por peso decrescente e seleciona iterativamente as disponíveis. \\

        \bottomrule
    \end{tabularx}
\end{table}


% ------------------------------------------------------------------------------------------------
% 4.3. Emparelhamento min-max
% ------------------------------------------------------------------------------------------------
\subsection{Problema de atribuição gargalo}

\subsubsection{Descrição do problema}

O problema de atribuição gargalo, também conhecido como \textbf{Problema de Emparelhamento Min-Max} ou \textbf{Bottleneck Assignment
problem}, é uma variação do problema de atribuição tradicional. O objetivo deste problema é encontrar, em um grafo bipartido ponderado, um emparelhamento de cardinalidade máxima (maior número possível de arestas) no qual o mínimo dos pesos das arestas do emparelhamento seja maximizado \cite{lawler}.

Por exemplo. Considere $N$ trabalhadores e $N$ estações de trabalho. $w_{ij}$ representa a eficiência do trabalhador i na estação de trabalho j. A eficiência total da produção é limitada pela eficiência do trabalhador menos eficiente.  O objetivo é atribuir trabalhadores às estações de trabalho de forma a maximizar a eficiência mínima entre todas as atribuições. \cite{lawler}.

\subsubsection{Soluções exatas}

\paragraph{Métodos de threshold}

Métodos de threshold funcionam alternando entre duas fases. A primeira fase consiste em selecionar um valor limite, chamado de \textbf{threshold}, descrito por $c\ast$. Na segunda fase, constrói-se um grafo bipartido com todas as arestas cujo peso é menor ou igual a $c\ast$. Em seguida, verifica-se se existe um emparelhamento perfeito nesse grafo. O menor valor de $c\ast$ que permite um emparelhamento perfeito é a solução ótima para o problema de atribuição gargalo \cite{burkard}.

Para selecionar o valor limite $c\ast$, podemos utilizar uma busca binária entre o menor e o maior peso das arestas do grafo original. A cada iteração, ajustamos o valor de $c\ast$ com base na existência ou não de um emparelhamento perfeito no grafo construído. Esse processo continua até que o valor ótimo seja encontrado. \citeonline{burkard}, que descreve esse método, utiliza um algoritmo baseado em matrizes para decidir se um emparelhamento é perfeito ou não, porém podemos utilizar qualquer algoritmo clássico de emparelhamento em grafos bipartidos, como os definidos na seção \ref{sec:algos_exatos_emp_card_max}, como por exemplo o algoritmo de Hopcroft-Karp. Também segundo o autor, o algoritmo tem uma complexidade de tempo de $O(V^{2.5} / \sqrt{\log V})$.

\citeonline{lawler} menciona um método de threshold para solução do problema de atribuição gargalo, porém a definição de threshold dele é levemente diferente, referindo-se ao valor dinâmico do rótulo de um vértice durante busca de caminhos. O algoritmo que ele descreve se assemelha mais aos algoritmos de caminhos aumentantes, que serão discutidos na seção \ref{sec:caminhos_aumentantes_min_max}.

\begin{algorithm}[H]
    \DontPrintSemicolon
    \SetAlgoLined

    % Main Algorithm Body
    \KwLet $C = (c_{ij})$ be a given $n \times n$ cost matrix;\;
    $c_0^* := \min_{ij}\{c_{ij}\}, c_1^* := \max_{ij}\{c_{ij}\}$\;
    
    \eIf{$c_0^* = c_1^*$}{
        $z := c_0^*$ \AlgComment{any permutation of $\{1, 2, \dots, n\}$ is optimal}
    }{
        \While{$C^* = \{c_{ij} : c_0^* < c_{ij} < c_1^*\} \neq \emptyset$}{
            \AlgComment{find the median, $c^*$, of $C^*$;}
            
            % The complex median formula
            $c^* := \min \{ c \in C^* : |\{ c_{ij} \in C^* : c_{ij} \le c \}| \ge |C^*|/2 \}$;\;
            
            Feasibility\_check($c^*, c_0^*, c_1^*$)\;
        }
        
        \If{Feasibility\_check has not yet been executed for $c_0^*$}{
             Feasibility\_check($c_0^*, c_0^*, c_1^*$)\;
        }
        
        find a perfect matching in the current bipartite graph $G[c_1^*]$\;
        $z := c_1^*$ \AlgComment{value of the optimal matching}
    }
    
    \bigskip % Add some vertical space between the main block and the procedure
    
    % The Procedure Definition
    \Fn{Feasibility\_check($c^*, c_0^*, c_1^*$)}{
        define the current graph $G[c^*]$\;
        \eIf{$G[c^*]$ contains a perfect matching}{
            $c_1^* := c^*$\;
        }{
            $c_0^* := c^*$\;
        }
    }
    
    \caption{Algoritmo Threshold para o LBAP \cite{burkard}}
    \label{alg:threshold}
\end{algorithm}

\paragraph{Método dual}

Assim como os métodos de threshold, o método dual também define um valor limite $c\ast$. No entanto, ao invés de 'chutar' um valor usando busca binária, este método define o valor inicial de $c\ast$ como sendo o menor valor possível que ele pode assumir. Olhando para a matriz de custos, esse valor é o máximo entre os mínimos de cada linha e cada coluna da matriz:

\begin{equation}
    c\ast = \max \left\{ \max_{i} (\min_{j} c_{ik}), \max_{j} (\min_{i} c_{kj}) \right\}
\end{equation}

Com o limite definido, verificamos se é possível encontrar um emparelhamento perfeito no grafo bipartido formado pelas arestas cujo peso é menor ou igual a $c\ast$, usando os mesmos algoritmos de antes. Se um emparelhamento perfeito for encontrado, então $c\ast$ é a solução ótima. Caso contrário, é necessário aumentar o valor de $c\ast$, e para isso nos baseamos no teorema de Kőnig, definido na seção \ref{p:teorema_konig}. 

Sabemos que todas as arestas $c_{ij} > c\ast$ estão conectadas a um conjunto de linhas $I$ e um conjunto de colunas $J$. Essas linhas e colunas "cobrem" todas as arestas do grafo atual. Para aumentar o valor de $c\ast$, precisamos adicionar uma nova aresta que não esteja coberta por esses conjuntos. Procuramos o menor peso entre as arestas não cobertas, ou seja, aquelas que não pertencem às linhas em $I$ nem às colunas em $J$. Esse valor se torna o novo $c\ast$:

\begin{equation}
    Novo\space c\ast = \min_{i \notin I, j \notin J} c_{ij}
\end{equation}

A partir daqui, o processo se repete. Com o novo valor de $c\ast$ e a nova aresta adicionada ao grafo, verificamos novamente a existência de um emparelhamento perfeito. Esse ciclo continua até que um emparelhamento perfeito seja encontrado, momento em que o valor atual de $c\ast$ será a solução ótima para o problema de atribuição gargalo.

\begin{algorithm}[H]
    \DontPrintSemicolon
    \SetAlgoLined

    % Line 1
    \KwLet $C = (c_{ij})$ be a given $n \times n$ cost matrix;\;
    
    % Line 2: Complex Max/Min assignment
    $c^* := \max_{k=1,2,\dots,n} (\min_{i=1,2,\dots,n} c_{ik}, \ \min_{j=1,2,\dots,n} c_{kj})$\;
    
    % Line 3
    $M := \emptyset$\;

    % Line 4: While Loop
    \While{$|M| < n$}{
        define the bipartite graph $G[c^*]$\;
        find a maximum matching $M$ in $G[c^*]$\;
        
        % Line 7: If statement
        \If{$|M| < n$}{
            % Line 8
            \KwLet $I \subseteq U$ and $J \subseteq V$ be vertex sets in a minimum vertex cover of $G[c^*]$\;
            
            % Line 9: Min assignment with 'not in' notation
            $c^* := \min_{i \notin I, \ j \notin J} c_{ij}$\;
        }
    }
    
    \caption{Dual\_LBAP: algoritmo dual para o LBAP \cite{burkard}}
    \label{alg:dual_lbap}
\end{algorithm}

O método dual não foi implementado no ambiente de testes deste trabalho. Para o problema de atribuição gargalo, foram comparados o algoritmo Threshold (Algoritmo~\ref{alg:threshold}) e o método de caminhos aumentantes (Algoritmo~\ref{alg:augmenting_lbap}), descrito a seguir.

\paragraph{Método de caminhos aumentantes}
\label{sec:caminhos_aumentantes_min_max}

Ao contrário do método dual, que ajusta um limiar global para todo o grafo, o método de caminhos aumentantes é construtivo. Ele inicia com um emparelhamento vazio e busca aumentar sua cardinalidade iterativamente, adicionando um par de vértices de cada vez. O objetivo central em cada passo é encontrar um caminho aumentante $P$ que conecte uma linha livre a uma coluna livre, de tal forma que a capacidade desse caminho seja mínima. A capacidade aqui é definida pelo peso da aresta mais cara contida no caminho:

\begin{equation}     
    cap(P) = \max_{(i,j) \in P} c_{ij} 
\end{equation}

Para encontrar tal caminho, utiliza-se uma adaptação do algoritmo de Dijkstra. Em vez de somar os pesos das arestas, mantemos para cada vértice um valor de rótulo que representa o "gargalo" mínimo necessário para alcançá-lo. Durante a busca, ao tentar estender o caminho de um nó $u$ (com rótulo $d_u$) para um vizinho $v$ através da aresta $(u,v)$, o custo não é somado, mas sim comparado. O novo custo potencial para atingir $v$ é o máximo entre o gargalo já existente até $u$ e o peso da nova aresta:

\begin{equation}     
    Novo\space Custo_v = \max(d_u, c_{uv}) 
\end{equation}

O algoritmo seleciona, a cada passo, o nó não visitado com o menor rótulo atual, garantindo uma exploração gulosa que minimiza o gargalo. Assim que uma coluna livre é alcançada por essa busca, o caminho aumentante é identificado e o emparelhamento é expandido, invertendo-se as arestas ao longo do caminho. Esse processo se repete (reiniciando as distâncias) até que um emparelhamento perfeito seja obtido. A solução do problema será determinada pela maior capacidade encontrada entre todos os caminhos aumentantes utilizados na construção.

\begin{algorithm}[H]
    \DontPrintSemicolon
    \SetAlgoLined

    \KwLet $G = (U, V; E)$ be a bipartite graph with edge lengths $c_{ij}$ for $[i, j] \in E$;\;
    find a lower bound $c^*$ for the optimum objective function value;\;
    \AlgComment{e.g., $c^* := \max(\max_{i \in U} \min_{j \in V} c_{ij}, \ \max_{j \in V} \min_{i \in U} c_{ij})$}\;
    
    find a maximum matching $M$ in the graph $G[c^*]$ having an edge $[i, j]$ \KwIff $c_{ij} \le c^*$;\;
    
    \eIf{$|M| = |U|$}{
        \KwStop \AlgComment{$M$ is an optimum matching of cost $c^*$}
    }{
        \KwLet $L$ be the set of unmatched vertices of $U$\;
    }

    % Continuation from the second image (The Main Loop)
    \While{$L$ is nonempty}{
        choose an arbitrary vertex $i \in L$;\;
        $L := L \setminus \{i\}$;\;
        Dijkstra($i$); \AlgComment{the procedure returns a path $P$ starting in $i$}\;
        
        \If{$P \neq \mathrm{nil}$}{
            $M := M \ominus P$;\;
            $c^* := \max(c^*, \ell(P))$\;
        }
    }
    \AlgComment{$M$ is a maximum matching with minimum cost $c^*$.}

    \caption{LBAP usando caminhos aumentantes \cite{burkard}}
    \label{alg:augmenting_lbap}
\end{algorithm}

\newpage

% ALGORITHM 6.4: DIJKSTRA AND LABEL
\begin{algorithm}[H]
    \DontPrintSemicolon
    \SetAlgoLined

    \Fn{Dijkstra($i$)}{
        label all vertices $j \in V$ by $(\alpha(j), \beta(j)) := (\infty, \mathrm{nil})$;\;
        $R := V$; \AlgComment{$R$ contains the unscanned vertices of $V$}\;
        $\overline{\alpha}(i) := c^*, P := \mathrm{nil}$;\;
        Label($i$);\;
        
        \While{$R \neq \emptyset$}{
            find a vertex $j_1 \in R$ with minimum $\alpha(j_1)$;\;
            \eIf{$\alpha(j_1) = \infty$}{
                $R := \emptyset$\;
            }{
                \eIf{$j_1$ is unmatched}{
                    find the path $P$ induced by the vertex sequence $(i, \dots, \overline{\beta}(\beta(j_1)), \beta(j_1), j_1)$;\;
                    $\ell(P) := \alpha(j_1)$;\;
                    $R := \emptyset$\;
                }{
                    let $[i_1, j_1]$ be the matching edge;\;
                    label $i_1$ with $(\overline{\alpha}(i_1), \overline{\beta}(i_1)) = (\alpha(j_1), j_1)$;\;
                    $R := R \setminus \{j_1\}$;\;
                    Label($i_1$);\;
                }
            }
        }
    }
    
    \bigskip
    
    \Fn{Label($i$)}{
        \For{each neighbor $j \in R$ of $i$}{
            \If{$\alpha(j) > \max(\overline{\alpha}(i), c_{ij})$}{
                $\alpha(j) := \max(\overline{\alpha}(i), c_{ij})$;\;
                $\beta(j) := i$\;
            }
        }
    }

    \caption{Dijkstra modificado para caminhos aumentantes em LBAP \cite{burkard}}
    \label{alg:dijkstra}
\end{algorithm}

\subsubsection{Heurísticas}

\paragraph{Heurística Construtiva Gulosa}

Diferente do problema de atribuição, onde o guloso simples funciona bem, no problema de gargalo o guloso deve ser adaptado para evitar criar "bicos de garrafa" muito altos no final do processo.

O algoritmo seleciona arestas iterativamente. Para cada vértice livre, ele verifica não apenas a aresta de menor peso incidente, mas também o "custo de oportunidade" de não escolhê-la.

Se não escolhermos a aresta $(u,v)$ de peso 10, a próxima melhor opção para $u$ tem peso 100? Se sim, a prioridade de $(u,v)$ aumenta drasticamente \cite{martello}.

% ### 3. Busca Local com Caminhos Curtos (Local Search)

% Uma maneira de melhorar um emparelhamento maximal (produzido pelo método guloso) é procurar por caminhos aumentantes curtos. Em grafos gerais, algoritmos exatos buscam caminhos aumentantes de *qualquer* tamanho (lidando com *blossoms*). A heurística limita essa busca a um tamanho fixo $k$.

% * **Descrição:**
%     1.  Gere um emparelhamento inicial $M$ (via Guloso).
%     2.  Procure por caminhos aumentantes de comprimento 3 (estrutura: aresta livre - aresta casada - aresta livre).
%     3.  Se encontrar um caminho $u - v - w - z$ onde $(v, w) \in M$, remova $(v, w)$ e adicione $(u, v)$ e $(w, z)$. O emparelhamento aumenta em +1.
% * **Complexidade:** Procurar caminhos de comprimento 3 pode ser feito em $O(E)$.
% * **Qualidade:** Apenas com caminhos de comprimento 3, a garantia de aproximação sobe para um fator melhor que 1/2.
% * **Fonte:** **Arkin, E. M., et al. (2006).** *Local search for the maximum cardinality matching problem*.

\begin{table}[htbp]
    \centering
    \caption{Algoritmos para problema de atribuição gargalo}
    \label{tab:atribuicao_gargalo_algos}
    \renewcommand{\arraystretch}{1.2} % Adds a little vertical padding
    
    % Structure: l = left align, c = center, X = auto-wrap paragraph
    \begin{tabularx}{\textwidth}{@{} l l l l X @{}}
        \toprule
        \textbf{Algoritmo} & \textbf{Complexidade} & \textbf{Tipo} & \textbf{Avaliado} & \textbf{Descrição} \\
        \midrule
        
        Threshold & $O(V^{2.5} / \sqrt{\log V})$ & Exato & Sim* & 
        Aplica busca binária sobre os pesos das arestas, testando a existência de emparelhamento perfeito no subgrafo restrito ao limiar. \\
        \addlinespace 
        
        Método Dual & $O(K \cdot E \sqrt{V})$ & Exato & Não & 
        Expande um limiar inferior iterativamente baseado no Teorema de König, adicionando apenas arestas necessárias para cobrir componentes. \\
        \addlinespace
        
        Caminhos aumentantes & $O(EV)$ & Exato & Sim & 
        Abordagem construtiva que busca caminhos aumentantes minimizando a aresta de peso máximo (gargalo) no trajeto (variação de Dijkstra). \\
        \addlinespace
        
        Construtiva Gulosa & $O(E \log E)$ & Heurística & Não & 
        Constrói a solução selecionando arestas de baixo peso, usando critérios locais ou de "arrependimento" para evitar gargalos futuros. \\
        
        \bottomrule
    \end{tabularx}
    \par\smallskip
    {\raggedright\footnotesize \textit{Avaliado}: incluído no ambiente de testes; a origem de cada implementação consta na Tabela~\ref{tab:proveniencia}. * A implementação avaliada testa cada limiar por caminhos aumentantes, com custo $O(n^3)$ por teste e $O(n^3 \log n)$ no total.\par}
\end{table}
% ------------------------------------------------------------------------------------------------
% 4.4. Emparelhamento estável (Stable Marriage Problem)
% ------------------------------------------------------------------------------------------------
\newpage
\subsection{Problema do emparelhamento estável}

\subsubsection{Descrição do problema}

O problema do emparelhamento estável, também conhecido como \textit{Stable Marriage Problem} (SMP), é um problema clássico na teoria dos grafos e na ciência da computação, proposto inicialmente por Gale e Shapley em 1962 \cite{kleinberg}. Ele envolve a combinação de dois conjuntos distintos de elementos e, apesar de semelhante aos problemas de emparelhamento em grafos bipartidos estudados anteriormente, possui características e objetivos diferentes.

Imagine dois conjuntos distintos de mesmo tamanho, $N$ homens e $N$ mulheres. Cada indivíduo em ambos os conjuntos tem uma lista de preferências ordenadas para todos os membros do outro conjunto. O objetivo do SMP é encontrar um emparelhamento entre homens e mulheres de tal forma que não existam dois indivíduos, um homem e uma mulher, que prefeririam estar juntos em vez de com seus parceiros atuais. Se tal par existir, o emparelhamento é considerado instável.

O problema pode ser resolvido utilizando o algoritmo de Gale-Shapley, que usa um processo iterativo de propostas e rejeições para garantir que o emparelhamento final seja estável.

% Variações do problema: Stable Roommates problem, Hospitais e Residentes, etc.
\subsubsection{Propriedades}

\begin{itemize}
    \item \textbf{Estabilidade}: Um emparelhamento é estável se não houver dois indivíduos que prefeririam estar juntos em vez de com seus parceiros atuais. A estabilidade é a característica central do problema \cite{lawler}.
    \item \textbf{Existência de solução}: Gale e Shapley provaram que sempre existe pelo menos um emparelhamento estável para qualquer conjunto de preferências \cite{ahuja}.
    \item \textbf{Otimalidade para quem propõe}: O emparelhamento resultante do algoritmo de Gale-Shapley é ótimo para o conjunto que faz as propostas (por exemplo, os homens), significando que cada indivíduo nesse conjunto recebe o melhor parceiro possível dentro dos emparelhamentos estáveis \cite{ahuja}.
    \item \textbf{Pessimismo para quem recebe propostas}: Por outro lado, o emparelhamento é o pior possível para o conjunto que recebe as propostas (por exemplo, as mulheres), significando que cada indivíduo nesse conjunto recebe o pior parceiro possível dentro dos emparelhamentos estáveis \cite{ahuja}.
    \item \textbf{Não-unicidade}: Pode haver múltiplos emparelhamentos estáveis para um dado conjunto de preferências, dependendo das listas de preferências dos indivíduos \cite{kleinberg}.
\end{itemize}

\subsubsection{Algoritmo de Gale-Shapley}

A solução clássica para o problema do emparelhamento estável é o algoritmo de Gale-Shapley, também conhecido como o algoritmo de casamento estável. O algoritmo funciona sob a lógica de propostas e rejeições, onde um dos conjuntos (por exemplo, os homens) faz propostas às mulheres com base em suas listas de preferências. 

O algoritmo pode ser descrito de forma simples da seguinte maneira:

Enquanto existir um homem $h$ que ainda não propôs a todas as mulheres em sua lista:

\begin{enumerate}
    \item O homem $h$ propõe à primeira mulher $m$ da sua lista de preferências a quem ele ainda não propôs.
    \item A decisão da mulher $m$ é tomada com base em suas preferências:
    \begin{itemize}
        \item Se $m$ não estiver comprometida, ela aceita a proposta de $h$. Eles ficam "noivos".
        \item Se $m$ já estiver comprometida com outro homem $h'$, ela compara $h$ e $h'$. Se $m$ preferir $h$ a $h'$, ela rejeita $h'$ (que fica livre) e aceita a proposta de $h$. Caso contrário, ela rejeita $h$ (que continua livre e tenta a próxima da sua lista).
    \end{itemize}
\end{enumerate}

O algoritmo continua até que todos os homens estejam comprometidos. O resultado final é um emparelhamento estável, onde não existem pares instáveis \cite{kleinberg}.

Segue um exemplo em pseudocódigo do algoritmo:

\begin{algorithm}[H]
    \DontPrintSemicolon
    \SetAlgoLined

    % Initial state
    Initially all $m \in M$ and $w \in W$ are free\;

    % Main Loop
    \While{there is a man $m$ who is free and hasn't proposed to every woman}{
        Choose such a man $m$\;
        Let $w$ be the highest-ranked woman in $m$'s preference list to whom $m$ has not yet proposed\;
        
        % Check if woman is free
        \eIf{$w$ is free}{
            $(m, w)$ become engaged\;
        }{
            % Else case: w is engaged to m'
            $w$ is currently engaged to $m'$\;
            
            \eIf{$w$ prefers $m'$ to $m$}{
                $m$ remains free\;
            }{
                % w prefers m to m'
                $(m, w)$ become engaged\;
                $m'$ becomes free\;
            }
        }
    }
    
    Return the set $S$ of engaged pairs\;

    \caption{Algoritmo de Gale-Shapley para o problema do emparelhamento estável \cite{kleinberg}}
    \label{alg:gale_shapley}
\end{algorithm}

Neste trabalho, o algoritmo de Gale-Shapley foi avaliado em três implementações, cuja origem consta na Tabela~\ref{tab:proveniencia}.

